\section{Summary}\label{sec:summary}
The implementation of RouteChasm emphasizes tight integration across all its systems.
Despite its large internal size, the framework maintains a clear architecture that supports extensibility
and long-term maintainability.

From creators point of view, the database layer implementation is intuitive, fast, and flexible.
It is by far the best designed system in RouteChasm.

The View architecture is also well-defined, although it bloats the project with deep directory structure.
This problem is planned to be addressed in the future.

Routing system went through many iterations and the latest one is the best, but it still has its flaws and quirks.
Implementation wise it is well-designed.
Usability wise it is not in a good spot.

The worst implementation is most likely the App class, which collects more and more technical debt with each new
feature.
The proper segmentation of App class has one of the highest priorities.

The Webpage Content Editor still uses outdated Editor which lacks crucial functionalities.
This component has the highest priority and in the near future it will have to be reimagined from scratch,
because the core architecture is not on par with the rest of the project.

AI system is implemented well, small tweaks are going to be needed.

All Framework Requirements are implemented and functional.
Some Non-Functional Requirements are loosely met.
Subjectively, the system performance is lacking [NFR3].
The performance is not noticeable, but 50ms execution time for a simple request is rather big toll.
RouteChasm is defined to be extensible and thus uses a lot of code from many files, which need to read, parsed, and
executed.
This decision adds delay to the execution process.
The framework is still use-able for small and medium size project where there may be only a few requests a second
during peak loads.
Some parts of the code do not follow the RouteChasms framework philosophy [NFR5].
This applies for legacy code or rushed feature code.
The culprits are going to be investigated in the near future and updated to follow the RouteChasm philosophy.

Overall I think that RouteChasm delivers on its goal well, it is robust maintainable extensible full-stack
solution for many web applications that may or may not use its AI functionalities.
